\section{Proofs for Deterministic and Statistical Recovery}
\label{app:recovery-proofs}

\begin{proof}[Proof of \cref{thm:joint-perturbation}]
Set
\[
  P=\Cpl(U,R),
  \qquad
  P_1=\Cpl(U,S),
  \qquad
  Q=\Cpl(V,S).
\]
At fixed \(U\), the completion section is isometric, so
\(\dai(P,P_1)=\dai(R,S)\).  Let \(U_t\) be the horizontal minimizing
Grassmann geodesic from \(U\) to \(V\), and keep \(S\) fixed in the parallel
frame.  The exact line element gives
\[
  \norm{\dot Q_t}_{Q_t,\AI}^2
  =2\tr(K_t\Xiop(S)K_t^\top)
  \le2\chi(S)^2\fro{K_t}^2.
\]
Therefore
\[
  \dai(P_1,Q)
  \le\sqrt2\chi(S)\int_0^1\fro{K_t}\,\dd t
  =\sqrt2\chi(S)\dgr(W,\widetilde W).
\]
The triangle inequality gives the claimed bound with \(\chi(S)\).  Reversing
the order, first rotating with \(R\) and then changing visible geometry at
fixed \(V\), gives the bound with \(\chi(R)\).  Taking the smaller proves
\eqref{eq:joint-perturbation}.
\end{proof}

\begin{proof}[Proof of \cref{thm:known-subspace-recovery}]
Let \(Y_i=R^{-1/2}X_i\).  Then
\[
  \E[Y_iY_i^\top]=I,
  \qquad
  \norm{Y_i}_2\le L,
\]
and
\[
  \widehat S=R^{-1/2}\Rhat R^{-1/2}
  =\frac1n\sum_{i=1}^nY_iY_i^\top.
\]
The \(Y_i\) are i.i.d.  For fixed unit \(u\), the variables
\(Z_i=(u^\top Y_i)^2\) satisfy \(0\le Z_i\le L^2\) and
\(\E Z_i=1\).  Hoeffding's inequality gives
\[
  \Pr\left\{
  \left|\frac1n\sum_iZ_i-1\right|\ge t
  \right\}
  \le2\exp\left(-\frac{2nt^2}{L^4}\right).
\]
Let \(\mathcal N\) be a \(1/4\)-net of the unit sphere in \(\R^m\), with
\(|\mathcal N|\le9^m\), and take \(t=\varepsilon/4\).  A union bound and the
sample-size assumption show that, with probability at least \(1-\delta\),
\[
  |u^\top(\widehat S-I)u|\le\varepsilon/4
  \qquad\forall u\in\mathcal N.
\]
For symmetric \(H\), a \(1/4\)-net obeys
\[
  \opnorm H\le2\sup_{u\in\mathcal N}|u^\top Hu|.
\]
Indeed, if unit \(x\) attains the operator norm and
\(\norm{x-u}\le1/4\), then
\(|x^\top Hx-u^\top Hu|\le\frac12\opnorm H\).  Hence
\(\opnorm{\widehat S-I}\le\varepsilon/2\).  For
\(0<\varepsilon\le1\),
\[
  e^{-\varepsilon}I
  \preceq(1-\varepsilon/2)I
  \preceq\widehat S
  \preceq(1+\varepsilon/2)I
  \preceq e^\varepsilon I.
\]
Congruence by \(R^{1/2}\) gives the Loewner sandwich.  Its generalized
log-eigenvalues have magnitude at most \(\varepsilon\), so
\(\dai(R,\Rhat)\le\sqrt m\,\varepsilon\).  Fixed-channel completion is
isometric, proving \eqref{eq:known-subspace-bound}.
\end{proof}

\begin{lemma}[Spectral-gap subspace perturbation]
\label{lem:spectral-subspace}
Let \(\Sigmahat=\Sigma+E\),
\(\opnorm E\le\delta<\gamma\), and let \(V\) span the leading
\(m\)-dimensional eigenspace of \(\Sigmahat\).  Then
\begin{equation}
  \opnorm{U_\perp^\top V}
  \le s_\delta=\frac{\delta}{\gamma-\delta}.
  \label{eq:sin-theta}
\end{equation}
If \(\delta<\gamma/2\), then
\begin{equation}
  \dgr(W,\widehat W)
  \le\sqrt m\,\arcsin(s_\delta),
  \label{eq:grassmann-stat}
\end{equation}
and after orthogonal Procrustes alignment,
\begin{equation}
  \opnorm{V-U}\le\sqrt2\,s_\delta.
  \label{eq:frame-stat}
\end{equation}
\end{lemma}

\begin{proof}
In the orthogonal basis \([U,U_\perp]\), write
\[
  \Sigma=\begin{pmatrix}A&0\\0&C\end{pmatrix},
  \qquad A\succeq\lambda_mI,
  \qquad C\preceq\lambda_{m+1}I.
\]
Let \(\widehat\Lambda\) contain the leading eigenvalues of \(\Sigmahat\).
Weyl's inequality gives
\(\widehat\Lambda\succeq(\lambda_m-\delta)I\).  With
\(Y=U_\perp^\top V\), projection of
\(\Sigmahat V=V\widehat\Lambda\) gives the Sylvester equation
\[
  Y\widehat\Lambda-CY=F,
  \qquad F=U_\perp^\top EV.
\]
The spectra are separated by at least \(\gamma-\delta\), and
\[
  Y=\int_0^\infty e^{tC}Fe^{-t\widehat\Lambda}\,\dd t.
\]
Thus
\[
  \opnorm Y
  \le\int_0^\infty e^{-t(\gamma-\delta)}\delta\,\dd t
  =\frac{\delta}{\gamma-\delta}.
\]
The singular values of \(U_\perp^\top V\) are the sines of the principal
angles, proving \eqref{eq:grassmann-stat}.  In principal frames,
\[
  \opnorm{V-U}
  =\max_i2\sin(\theta_i/2)
  \le\sqrt2\max_i\sin\theta_i
  \le\sqrt2s_\delta,
\]
which proves \eqref{eq:frame-stat}.
\end{proof}

\begin{proof}[Proof of \cref{thm:unknown-subspace-recovery}]
Align \(V\) to \(U\) as in \cref{lem:spectral-subspace}, and conjugate
\(\Rhat\) by the same orthogonal alignment.  Denote the aligned matrix by
\(S\); the completed estimator is unchanged.  Then
\begin{align*}
  \opnorm{S-R}
  &\le\opnorm{V^\top(\Sigmahat-\Sigma)V}
  +\opnorm{V^\top\Sigma V-U^\top\Sigma U}\\
  &\le\delta+2\lambda_1\opnorm{V-U}\\
  &\le\delta+2\sqrt2\lambda_1s_\delta
  =\varepsilon_{\mathrm{vis}}.
\end{align*}
Since \(R\succeq\lambda_mI\),
\[
  \opnorm{R^{-1/2}(S-R)R^{-1/2}}
  \le\eta_\delta<1.
\]
Therefore
\[
  (1-\eta_\delta)R\preceq S\preceq(1+\eta_\delta)R,
\]
and
\begin{equation}
  \dai(R,S)\le\sqrt m[-\log(1-\eta_\delta)].
  \label{eq:aligned-visible-bound}
\end{equation}
Weyl's inequality places the eigenvalues of \(S\) in
\([\lambda_m-\delta,\lambda_1+\delta]\), so
\(\chi(S)\le\omega_\delta\).  Apply
\cref{thm:joint-perturbation,lem:spectral-subspace} and
\eqref{eq:aligned-visible-bound}:
\begin{align*}
  \dai(P,\Phat)
  &\le\dai(R,S)+\sqrt2\chi(S)\dgr(W,\widehat W)\\
  &\le\sqrt m[-\log(1-\eta_\delta)]
  +\sqrt2\omega_\delta\sqrt m\arcsin(s_\delta).
\end{align*}
This is \eqref{eq:unknown-subspace-bound}.
\end{proof}

\begin{lemma}[Ambient second-moment concentration]
\label{lem:ambient-concentration}
Under the assumptions of \cref{cor:unknown-subspace-sample}, with probability
at least \(1-\alpha\),
\[
  \opnorm{\Sigmahat-\Sigma}\le\delta_n(\alpha).
\]
\end{lemma}

\begin{proof}
Let \(\mathcal N\) be a \(1/4\)-net of the Euclidean unit sphere with
\(|\mathcal N|\le9^d\).  For fixed \(u\in\mathcal N\), the i.i.d. variables
\(Z_i=(u^\top X_i)^2\) lie in \([0,L^2]\).  Hoeffding and a union bound give
\[
  \Pr\left\{
  \sup_{u\in\mathcal N}|u^\top(\Sigmahat-\Sigma)u|\ge t
  \right\}
  \le2\cdot9^d\exp\left(-\frac{2nt^2}{L^4}\right).
\]
The net estimate used above gives
\(\opnorm{\Sigmahat-\Sigma}\le2t\).  Taking
\(t=\delta_n(\alpha)/2\) makes the failure probability at most \(\alpha\).
\end{proof}

\begin{proof}[Proof of \cref{cor:unknown-subspace-sample}]
On the event from \cref{lem:ambient-concentration}, all assumptions of
\cref{thm:unknown-subspace-recovery} hold with
\(\delta=\delta_n(\alpha)\).  Apply that theorem.  For the explicit sufficient
condition, \(\delta\le\delta_\star\le\gamma/4\) gives
\[
  s_\delta=\frac{\delta}{\gamma-\delta}
  \le\frac{4\delta}{3\gamma}.
\]
Consequently
\[
  \varepsilon_{\mathrm{vis}}
  \le\delta\left(1+\frac{8\sqrt2\lambda_1}{3\gamma}\right)
  =c_\gamma\delta
  \le\frac{\lambda_m}{2},
  \qquad
  \eta_\delta\le\frac12.
\]
Also \(\delta\le\gamma/4<\gamma/2\) and
\(\delta\le\lambda_m/(2c_\gamma)<\lambda_m\), so every hypothesis of
\cref{thm:unknown-subspace-recovery} is satisfied.
\end{proof}
